% ===== 原文 pp.28–31 =====
% 本块从 p.28 的 §5 标题起（p.28 上半页仍是 tr-04 的 §4.3.3，此处不重复）
% [原文 p.28]

\section{两个对应关系及其解释}

关键之处到了：我们要把第 3 节引入的、与带穿刺点的 Riemann 曲面相联系的那一类四维
$\cN=2$ 理论，放到第 4 节引入的两个超对称背景上，计算它们的配分函数。我们会发现，
它们由第 2 节引入的两个二维理论给出。

\subsection{$S^1\times S^3$ 与二维杨--米尔斯理论}

我们先考虑 $S^1\times S^3$ 上的配分函数。取一个三穿刺球面。在第 3 节中，我们给它联系了
$\SU(2)_1\times\SU(2)_2\times\SU(2)_3$ 的 $2\otimes2\otimes2$ 中的一个半超多重态。它的%
% [原文 p.29]
配分函数 $Z(a_1,a_2,a_3)$ 在 $S^1\times S^3$ 上——当 $\SU(2)_i$ 的和乐为
$(a_i,1/a_i)\in\SU(2)_i$ 时——由 (4.9) 给出为
\begin{equation}
Z(a_1,a_2,a_3)=\prod_{\pm\pm\pm}\prod_{n\ge0}\frac{1}{1-q^{n+1/2}a_1^{\pm1}a_2^{\pm1}a_3^{\pm1}}.
\tag{5.1}
\end{equation}

% [Tutor Note: p.29, (5.1) 的 $\prod_{\pm\pm\pm}$ 跑遍 $2\otimes2\otimes2$ 的全部 8 个权——半超多重态
%  在超共形指标里的贡献就是对表示的全部权取乘积（这里 $2\otimes2\otimes2$ 是伪实表示）。
%  自检办法：把 (5.2) 右端按 $q^{1/2}$ 展开，其首项为
%  $q^{1/2}\chi_2(a_1)\chi_2(a_2)\chi_2(a_3)=q^{1/2}\sum_{8\text{个权}}a^w$，与 (5.1) 的
%  $\prod_{\pm\pm\pm}$ 首项完全一致；若 (5.1) 只取 4 项，首项会少一半，与 (5.2) 对不上。
%  注意 p.26 关于「半超多重态只取一半的权」的规则只适用于 $S^4$ 上的单圈行列式
%  （见 (5.13) 的 $\prod_{\pm\pm}$），不适用于 $S^1\times S^3$ 的超共形指标。]

它恰好还有一个无穷和形式的等价表达式
\begin{equation}
Z(a_1,a_2,a_3)=\frac{K(a_1)K(a_2)K(a_3)}{K_0}\sum_{n\ge0}
\frac{\chi_n(a_1)\chi_n(a_2)\chi_n(a_3)}{\chi_n(q^{1/2})},
\tag{5.2}
\end{equation}
其中
\begin{equation}
K(a)^{-1}=\prod_{n\ge0}\left[(1-q^{n+1})\prod_{\pm}(1-q^{n+1}a^{\pm2})\right]
\tag{5.3}
\end{equation}
就是 (4.11) 中引入的同一个函数，
\begin{equation}
K_0^{-1}=\prod_{n\ge0}(1-q^{2+n}),
\tag{5.4}
\end{equation}
而
\begin{equation}
\chi_n(a)=a^{n-1}+a^{n-3}+\cdots+a^{3-n}+a^{1-n}
\tag{5.5}
\end{equation}
是 $(a,1/a)\in\SU(2)$ 在 $\SU(2)$ 的 $n$ 维不可约表示中的特征标。

% [Tutor Note: p.29, 紧接 (5.5) 的这句话里、以及下面「$K(0)$」处，原文有两处瑕疵
%  （均已用 800 dpi 裁图逐字复核）。其一，「introduced in Sec. 2.30」所指的小节并不存在：
%  §2 只有 2.1、2.2 及其 $x.y.z$ 子节；从上下文看这里应指 §2.1.4（q-形变，原文 pp.8--9），
%  三孔球面振幅与式 (2.30) 都在该小节引入。其二，「prefactor involving $K(a_i)$ and
%  $K(0)$」中的 $K(0)$ 应为 $K_0$：由 (5.4) 定义、且在 (5.2)、(5.11) 中用的都是 $K_0$；
%  而按 (5.3)，$K(a)$ 在 $a=0$ 处因 $a^{\pm2}$ 而发散，$K(0)$ 本身没有定义。正文照原文排印。]

至此，读者不妨试着证明等式 (5.2)，或者至少把两边展开到 $O(q^3)$ 量级来核对一下。

我们看到，这个无穷和表达式等于第 2.30 节引入的二维 $q$-形变 $\SU(2)$ 杨--米尔斯理论中
三孔球面振幅在零面积极限下的值，只差一个与 $K(a_i)$ 和 $K(0)$ 有关的预因子。

下面考虑四穿刺球面。相应的四维理论，是把 $\SU(2)_1\times\SU(2)_2\times\SU(2)$ 与
$\SU(2)\times\SU(2)_3\times\SU(2)_4$ 的两个半超多重态，用动力学的 $\SU(2)$ 规范多重态
耦合起来得到的。相应地，$S^1\times S^3$ 上的配分函数由公式 (4.10) 给出：
\begin{align}
Z(a_1,a_2;a_3,a_4)&=\frac{1}{2}\oint\frac{dz}{2\pi iz}(1-z^2)\left(1-\frac{1}{z^2}\right)
K(z)^{-2}Z(a_1,a_2,z)Z(z,a_3,a_4)
\tag{5.6}\\
&=\frac{1}{2}\oint\frac{dz}{2\pi iz}(1-z^2)\left(1-\frac{1}{z^2}\right)K(z)^{-2}\notag\\
&\quad\times\prod_{\pm\pm\pm}\prod_{n\ge0}\frac{1}{1-q^{n+1/2}a_1^{\pm1}a_2^{\pm1}z^{\pm1}}
\prod_{\pm\pm\pm}\prod_{n\ge0}\frac{1}{1-q^{n+1/2}z^{\pm1}a_3^{\pm1}a_4^{\pm1}}
\tag{5.7}
\end{align}
这里我们用了半超多重态贡献的无穷乘积形式 (5.1)。从这一形式看不出该表达式在变量
$a_{1,2,3,4}$ 的置换下是对称的。

% [原文 p.30]

若改用无穷和形式 (5.2) 代入，则得到
\begin{align}
Z(a_1,a_2;a_3,a_4)&=\frac{1}{2}\oint\frac{dz}{2\pi iz}(1-z^2)\left(1-\frac{1}{z^2}\right)
\frac{K(a_1)K(a_2)K(a_3)K(a_4)}{K_0^2}\notag\\
&\quad\times\sum_{n\ge0}\frac{\chi_n(a_1)\chi_n(a_2)\chi_n(z)}{\chi_n(q^{1/2})}
\sum_{m\ge0}\frac{\chi_m(z)\chi_m(a_2)\chi_n(a_3)}{\chi_m(q^{1/2})}
\tag{5.8}\\
&=\frac{K(a_1)K(a_2)K(a_3)K(a_4)}{K_0^2}\sum_{n\ge0}
\frac{\chi_n(a_1)\chi_n(a_2)\chi_m(a_3)\chi_n(a_4)}{\chi_n(q^{1/2})^2}.
\tag{5.9}
\end{align}

% [Tutor Note: p.30, (5.8)、(5.9) 的指标有两处笔误（已用 1400 dpi 裁图逐字复核）。
%  (5.8) 第二个求和的分子原文印作 $\chi_m(z)\chi_m(a_2)\chi_n(a_3)$，但它是三基表示振幅
%  $Z(z,a_3,a_4)$ 按 (5.2) 展开的因子，按 (5.6) 的结构应为
%  $\chi_m(z)\chi_m(a_3)\chi_m(a_4)$；(5.9) 分子中第三个因子原文印作 $\chi_m(a_3)$，
%  用 (5.10) 正交性把 $z$ 积掉后应为 $\chi_n(a_3)$（其余三个因子下标都是 $n$）。
%  正文照原文排印，此处只作记录；读者若自己重算 (5.7) 会直接得到修正后的形式。]

这里我们用了 $\SU(2)$ 特征标在自然测度下的正交性
\begin{equation}
\frac{1}{2}\oint\frac{dz}{2\pi iz}(1-z^2)\left(1-\frac{1}{z^2}\right)\chi_n(z)\chi_m(z)=\delta_{mn}
\tag{5.10}
\end{equation}
它不过是 (2.6) 的一个特例。因子 $(1-z^2)(1-z^{-2})$ 是把群变量 $z\in\SU(2)$ 限制到其
Cartan 环面上时引入的测度因子。

% [Tutor Note: p.30, (5.10) 中的 $1/2$ 来自 Weyl 群 $\Z_2$：$\SU(2)$ 的群积分限制到
%  Cartan 环面上后成为 $\frac{1}{|W|}\oint\frac{dz}{2\pi iz}(1-z^2)(1-z^{-2})$，其中
%  $(1-z^2)(1-z^{-2})$ 是 Weyl 分母的平方（差一个符号）。用 (5.5) 的 $\chi_n$ 代入即可直接验证 $\delta_{mn}$；
%  这一步与 (2.6) 是同一个测度，只是换了记号。]

我们发现，配分函数 (5.9) 又一次等于 $q$-形变 $\SU(2)$ 杨--米尔斯理论中四孔球面的振幅，
只不过差一个与 $K(a)$ 和 $K_0$ 有关的预因子，并且面积 $A$ 需要取为零。

这一计算可以容易地推广到任意的 $\SU(2)$ 型 class S 理论。在四维一侧，我们考虑与亏格为
$g$、有 $n$ 个穿刺点的 Riemann 曲面相联系的理论在 $S^1\times S^3$ 上的配分函数，其中对第
$i$ 个穿刺点的味对称性取和乐 $(a_i,1/a_i)\in\SU(2)_i$。在二维一侧，我们在同一个 Riemann
曲面上取 $q$-形变 $\SU(2)$ 杨--米尔斯理论，把穿刺点看作孔，并令第 $i$ 个孔周围的和乐为
$(a_i,1/a_i)\in\SU(2)$。于是有一般的关系式
\begin{equation}
Z_{4\mathrm{d},g,n}(a_i)=\frac{\prod_iK(a_i)}{K_0^{2g-2+n}}Z_{2\mathrm{d},g,n,A=0}(a_i).
\tag{5.11}
\end{equation}

每个穿刺点带来的因子 $K(a)/K_0$ 可以吸收进二维一侧「引入孔的算符」的重定义之中。
类似地，在二维非引力理论中，我们总可以加上如下形式的局部抵消项
\begin{equation}
S_{2\mathrm{d}}\to S_{2\mathrm{d}}+c\int\sqrt{g}\,R_g=S_{2\mathrm{d}}+c(2-2g)
\tag{5.12}
\end{equation}
其中 $R_g$ 是二维度规的曲率标量。这可以吸收因子 $K_0^{2-2g}$。因此我们得出结论：与任意
带穿刺点的 Riemann 曲面相联系的 $\SU(2)$ 型 class S 理论，其在 $S^1\times S^3$ 上的配分
函数总等于在同一带穿刺 Riemann 曲面上考虑的 $q$-形变 $\SU(2)$ 杨--米尔斯理论的配分函数。
这就是最早在 [3] 中发现的对应关系。

% [Tutor Note: p.30, (5.11) 到 (5.12) 的计数：$K_0^{2g-2+n}$ 在分母上，$n$ 个穿刺点各吸收一个
%  $K(a_i)/K_0$ 后，剩下的是 $K_0^{-(2g-2)}=K_0^{2-2g}$，与正文「absorb the factor $K_0^{2-2g}$」
%  一致。(5.12) 的拓扑项在欧氏路径积分里给出的因子是 $e^{-c(2-2g)}$；要让它恰好等于 $K_0^{2-2g}$，
%  需 $e^{-c}=K_0$，即 $c=-\log K_0$——它是 $q$ 的函数，一般并不等于 $-1$。原文只说「这可以吸收
%  该因子」，没有指定 $c$ 的数值。易错点：$(2-2g)$ 与 $(2g-2)$ 互为相反数，写成哪一边取决于
%  因子在分子还是分母，不要按「乘以 $K_0^{c(2-2g)}$」去类比。]

% [原文 p.31]

\subsection{$S^4$ 与 Liouville 理论}

下面考虑 $S^4$ 上的配分函数，它将是超引力背景中的参数 $\epsilon_{1,2}$、各质量 $m_i$
以及复化规范耦合 $\tau_i$ 的函数。注意 $\epsilon_{1,2}$ 与 $m_i$ 的量纲都是 1。我们通过
要求 $\epsilon_1\epsilon_2=1$ 来固定标度。

取一个三穿刺球面。第 3 节中与它相联系的四维理论，是
$\SU(2)_1\times\SU(2)_2\times\SU(2)_3$ 的 $2\otimes2\otimes2$ 中的一个半超多重态。当
$\SU(2)_i$ 的质量参数为 $(m_i/2,-m_i/2)\in\SU(2)_i$ 时，它在 $S^4$ 上的配分函数
$Z^{\text{trifund}}_{S^4}(m_1,m_2,m_3)$ 由 $|Z_{\text{n.h.}}(m_1,m_2,m_3)|^2$ 给出，如 (4.15)
所示。这里来自北半球的贡献 $Z^{\text{trifund}}_{\text{n.h.}}(m_1,m_2,m_3)$ 由 (4.19) 给出：
\begin{equation}
Z^{\text{trifund}}_{\text{n.h.}}(m_1,m_2,m_3)=\prod_{\pm\pm}
\Gamma_2\left(\frac{m_1\pm m_2\pm m_3}{2}+\frac{\epsilon_1+\epsilon_2}{2}\,\Big|
\,\epsilon_1,\epsilon_2\right)
\tag{5.13}
\end{equation}
因此我们有
\begin{equation}
Z^{\text{trifund}}_{S^4}(m_1,m_2,m_3)=\prod_{\pm\pm\pm}
\Gamma_2\left(\frac{\pm m_1\pm m_2\pm m_3}{2}+\frac{\epsilon_1+\epsilon_2}{2}\,\Big|
\,\epsilon_1,\epsilon_2\right).
\tag{5.14}
\end{equation}
注意这里用到了我们的约定：$m_i$ 全都是纯虚的。

% [Tutor Note: p.31, (5.13) 是 $\prod_{\pm\pm}$（4 项），(5.14) 却是 $\prod_{\pm\pm\pm}$（8 项），
%  这个变化原文一句话带过，补出来是：按正文 (5.14) 等于 $|Z_{\text{n.h.}}|^2$，即把 (5.13)
%  与其复共轭相乘；由 (2.50) 的 $\Gamma_2(x^*|\epsilon_1,\epsilon_2)=\Gamma_2(x|\epsilon_1,
%  \epsilon_2)^*$、本页约定的 $m_i$ 纯虚（故 $m_i^*=-m_i$）以及 $\epsilon_{1,2}$ 实，
%  复共轭恰好把辐角中的 $m_1\pm m_2\pm m_3$ 翻成 $-(m_1\pm m_2\pm m_3)$，于是 4 项乘上
%  4 个共轭项，合起来正是 8 项 $=\prod_{\pm\pm\pm}$。]

我们看到，在如下识别下，它等于 Liouville 理论三点函数 (2.46) 的分母：
\begin{equation}
\alpha_i=m_i+\frac{b+1/b}{2},\qquad(\epsilon_1,\epsilon_2)=\left(b,\frac{1}{b}\right),
\tag{5.15}
\end{equation}
这里已经代入了函数 $\Upsilon$ 的定义 (2.48)。分子中的大部分因子也可以通过把如下二维算符
\begin{equation}
V_{\alpha_i}:=\Upsilon(2\alpha_i)e^{2\alpha_i\phi}
\tag{5.16}
\end{equation}
等同于 class S 构造中质量为 $m_i$ 的穿刺点来交代。这样一来，Liouville 三点函数中大部分
依赖于质量的项就落在三基表示半超多重态的配分函数里。把这个写成一个等式：
\begin{equation}
Z^{\text{trifund}}_{S^4}(m_1,m_2,m_3)=\langle V_{\alpha_1}(\infty)V_{\alpha_2}(1)
V_{\alpha_3}(0)\rangle_{\Liouville}
\tag{5.17}
\end{equation}
其中等号成立到乘以一个与 $m_i$ 无关的函数为止。

现在考虑四穿刺球面。相应的四维理论这样得到：取两个三基表示，一个属于
$\SU(2)_1\times\SU(2)_2\times\SU(2)$，另一个属于
$\SU(2)\times\SU(2)_3\times\SU(3)_4$，再把它耦合到 $\SU(2)$ 矢量多重态上。按 (4.15)，
$S^4$ 配分函数为

% [Tutor Note: p.31, 上面这一段把第二个三基表示的第三个因子印作
%  $\SU(2)\times\SU(2)_3\times\SU(3)_4$，即写成 $\SU(3)_4$。这是笔误：紧随其后的下文
%  （原文 p.32）以及三基表示的定义都要求三个因子都是 $\SU(2)$，应为
%  $\SU(2)\times\SU(2)_3\times\SU(2)_4$。（已用 800 dpi 裁图复核，原文确实印的是
%  $\SU(3)_4$。）正文照原文排印。]

\begin{align}
Z_{S^4}(m_1,m_2;m_3,m_4;\tau)&=\frac{1}{2}\int da\,|q|^{2a^2}
Z^{\text{trifund}}_{S^4}(m_1,m_2,a)Z^{\text{trifund}}_{S^4}(a,m_3,m_4)\notag\\
&\quad\times Z^{\text{gauge,one-loop}}_{S^4}(a)\left|Z^{\text{instanton}}_{\text{n.p.}}
(a,m_i,q)\right|^2.
\tag{5.18}
\end{align}

